% GENERATED FILE. Edit canonical lessons, then run npm run generate:books.
% canonical-source-hash: fnv1a64:896fcc7fcb531f83
% canonical-lessons: BN-C335-osvikar-kora, BN-C335-somalochona-kora, BN-C335-drishtibhonggi, BN-C335-udahoron, BN-C335-prosonggo

\chapter{Granting and Denying}
\label{ch:bn-granting-and-denying}
\hlchaptermodality{335}

\begin{chapteropening}
\textbf{By the end of this chapter:} \emph{I can deny a claim, criticise fairly and give an example.}

Complete the last lesson of chapter 335: I can deny a claim, criticise fairly and give an example.
\end{chapteropening}

\section[\texorpdfstring{\bn{অস্বীকার} \bn{করা} (ôsvīkār kôrā)}{অস্বীকার করা (ôsvīkār kôrā)}]{\bn{অস্বীকার} \bn{করা} (ôsvīkār kôrā) — to deny}

\label{lesson:BN-C335-osvikar-kora}



\begin{quote}
Before the new one: say the Bengali for to oppose, then the Bengali for to disregard.
\end{quote}

\subsection*{You'll want to know: \bn{অস্বীকার} \bn{করা}}
\textbf{\bn{অস্বীকার} \bn{করা}} — \emph{ôsvīkār kôrā} — \textquotedblleft{}to deny\textquotedblright{}.

\begin{quote}
\bn{চোর} \bn{সব} \bn{অভিযোগ} \bn{অস্বীকার} \bn{করল।}
\end{quote}

\begin{quote}
\emph{chor sôb ôbhijog ôsvīkār kôrlô.}
\end{quote}

\begin{quote}
The thief denied every charge.
\end{quote}

Past (he/she …): \bn{তিনি} \bn{কথাটা} \bn{অস্বীকার} \bn{করলেন}, \emph{tini kôthāṭā ôsvīkār kôrlen}.

\begin{cousinweb}
\textbf{\bn{অস্বীকার} \bn{করা}} (\emph{ôsvīkār kôrā}) is Sanskrit \emph{a-} 'not' with \emph{svīkāra}, assent or acceptance (Monier-Williams), and \bn{করা}: to refuse to accept that something is true.
\end{cousinweb}

\begin{grammarlens}[title={What you've built}]
One more word for giving reasons and opinions.
\end{grammarlens}

\subsection*{Your turn}
\begin{itemize}
\raggedright
  \item \emph{Say it:} \emph{ôsvīkār kôrā}
  \item \emph{Say it:} \emph{ôsvīkār kôrā}, once more
  \item \emph{Say it:} \emph{chor sôb ôbhijog ôsvīkār kôrlô.} — The thief denied every charge.
  \item \emph{Recall:} say the Bengali for to oppose, then the Bengali for to disregard, then say \emph{ôsvīkār kôrā} again
\end{itemize}

\begin{tcolorbox}[breakable,colback=teal!4,colframe=teal!35!black,title={Before you move on}]
What does \bn{অস্বীকার} \bn{করা} mean? (\textquotedblleft{}To deny\textquotedblright{}.) Say it once more.
\end{tcolorbox}
\section[\texorpdfstring{\bn{সমালোচনা} \bn{করা} (sômālochônā kôrā)}{সমালোচনা করা (sômālochônā kôrā)}]{\bn{সমালোচনা} \bn{করা} (sômālochônā kôrā) — to criticise}

\label{lesson:BN-C335-somalochona-kora}



\begin{quote}
Before the new one: say the Bengali for to disregard, then the Bengali for to deny.
\end{quote}

\subsection*{You'll want to know: \bn{সমালোচনা} \bn{করা}}
\textbf{\bn{সমালোচনা} \bn{করা}} — \emph{sômālochônā kôrā} — \textquotedblleft{}to criticise\textquotedblright{}.

\begin{quote}
\bn{সাংবাদিকরা} \bn{সরকারের} \bn{সমালোচনা} \bn{করলেন।}
\end{quote}

\begin{quote}
\emph{sāngbādikrā sôrkārer sômālochônā kôrlen.}
\end{quote}

\begin{quote}
The journalists criticised the government.
\end{quote}

Past (he/she …): \bn{তিনি} \bn{বইটার} \bn{সমালোচনা} \bn{করলেন}, \emph{tini bôiṭār sômālochônā kôrlen}.

\begin{cousinweb}
\textbf{\bn{সমালোচনা} \bn{করা}} (\emph{sômālochônā kôrā}) is built on Sanskrit \emph{ālocana}, considering or reflecting (Monier-Williams), with \emph{sam-} 'thoroughly'. The sense 'criticism', a close look that finds faults, is a later Indic use. What is criticised takes -\bn{এর}.
\end{cousinweb}

\begin{grammarlens}[title={What you've built}]
One more word for giving reasons and opinions.
\end{grammarlens}

\subsection*{Your turn}
\begin{itemize}
\raggedright
  \item \emph{Say it:} \emph{sômālochônā kôrā}
  \item \emph{Say it:} \emph{sômālochônā kôrā}, once more
  \item \emph{Say it:} \emph{sāngbādikrā sôrkārer sômālochônā kôrlen.} — The journalists criticised the government.
  \item \emph{Recall:} say the Bengali for to disregard, then the Bengali for to deny, then say \emph{sômālochônā kôrā} again
\end{itemize}

\begin{tcolorbox}[breakable,colback=teal!4,colframe=teal!35!black,title={Before you move on}]
What does \bn{সমালোচনা} \bn{করা} mean? (\textquotedblleft{}To criticise\textquotedblright{}.) Say it once more.
\end{tcolorbox}
\section[\texorpdfstring{\bn{দৃষ্টিভঙ্গি} (drishṭibhônggi)}{দৃষ্টিভঙ্গি (drishṭibhônggi)}]{\bn{দৃষ্টিভঙ্গি} (drishṭibhônggi) — a point of view; an outlook}

\label{lesson:BN-C335-drishtibhonggi}



\begin{quote}
Before the new one: say the Bengali for to deny, then the Bengali for to criticise.
\end{quote}

\subsection*{You'll want to know: \bn{দৃষ্টিভঙ্গি}}
\textbf{\bn{দৃষ্টিভঙ্গি}} — \emph{drishṭibhônggi} — \textquotedblleft{}a point of view; an outlook\textquotedblright{}.

\begin{quote}
\bn{তার} \bn{দৃষ্টিভঙ্গি} \bn{আমার} \bn{থেকে} \bn{আলাদা।}
\end{quote}

\begin{quote}
\emph{tār drishṭibhônggi āmār theke ālādā.}
\end{quote}

\begin{quote}
Her outlook is different from mine.
\end{quote}

\begin{cousinweb}
\textbf{\bn{দৃষ্টিভঙ্গি}} (\emph{drishṭibhônggi}) is Sanskrit \emph{dṛṣṭi}, sight, with \emph{bhaṅgi}, a bend or a manner (Monier-Williams): the angle from which you look. The compound is a modern one.
\end{cousinweb}

\begin{grammarlens}[title={What you've built}]
One more word for giving reasons and opinions.
\end{grammarlens}

\subsection*{Your turn}
\begin{itemize}
\raggedright
  \item \emph{Say it:} \emph{drishṭibhônggi}
  \item \emph{Say it:} \emph{drishṭibhônggi}, once more
  \item \emph{Say it:} \emph{tār drishṭibhônggi āmār theke ālādā.} — Her outlook is different from mine.
  \item \emph{Recall:} say the Bengali for to deny, then the Bengali for to criticise, then say \emph{drishṭibhônggi} again
\end{itemize}

\begin{tcolorbox}[breakable,colback=teal!4,colframe=teal!35!black,title={Before you move on}]
What does \bn{দৃষ্টিভঙ্গি} mean? (\textquotedblleft{}A point of view; an outlook\textquotedblright{}.) Say it once more.
\end{tcolorbox}
\section[\texorpdfstring{\bn{উদাহরণ} (udāhôrôn)}{উদাহরণ (udāhôrôn)}]{\bn{উদাহরণ} (udāhôrôn) — an example}

\label{lesson:BN-C335-udahoron}



\begin{quote}
Before the new one: say the Bengali for to criticise, then the Bengali for an outlook.
\end{quote}

\subsection*{You'll want to know: \bn{উদাহরণ}}
\textbf{\bn{উদাহরণ}} — \emph{udāhôrôn} — \textquotedblleft{}an example\textquotedblright{}.

\begin{quote}
\bn{একটা} \bn{উদাহরণ} \bn{দাও।}
\end{quote}

\begin{quote}
\emph{ekṭā udāhôrôn dāo.}
\end{quote}

\begin{quote}
Give an example.
\end{quote}

\begin{cousinweb}
\textbf{\bn{উদাহরণ}} (\emph{udāhôrôn}) is Sanskrit \emph{udāharaṇa}, an example or illustration (Monier-Williams). You give one with \bn{দেওয়া}: \bn{উদাহরণ} \bn{দাও}.
\end{cousinweb}

\begin{grammarlens}[title={What you've built}]
One more word for giving reasons and opinions.
\end{grammarlens}

\subsection*{Your turn}
\begin{itemize}
\raggedright
  \item \emph{Say it:} \emph{udāhôrôn}
  \item \emph{Say it:} \emph{udāhôrôn}, once more
  \item \emph{Say it:} \emph{ekṭā udāhôrôn dāo.} — Give an example.
  \item \emph{Recall:} say the Bengali for to criticise, then the Bengali for an outlook, then say \emph{udāhôrôn} again
\end{itemize}

\begin{tcolorbox}[breakable,colback=teal!4,colframe=teal!35!black,title={Before you move on}]
What does \bn{উদাহরণ} mean? (\textquotedblleft{}An example\textquotedblright{}.) Say it once more.
\end{tcolorbox}
\section[\texorpdfstring{\bn{প্রসঙ্গ} (prôsônggô)}{প্রসঙ্গ (prôsônggô)}]{\bn{প্রসঙ্গ} (prôsônggô) — a topic of talk; a context}

\label{lesson:BN-C335-prosonggo}



\begin{quote}
Before the new one: say the Bengali for an outlook, then the Bengali for an example.
\end{quote}

\subsection*{You'll want to know: \bn{প্রসঙ্গ}}
\textbf{\bn{প্রসঙ্গ}} — \emph{prôsônggô} — \textquotedblleft{}a topic of talk; a context\textquotedblright{}.

\begin{quote}
\bn{এই} \bn{প্রসঙ্গে} \bn{আমার} \bn{একটা} \bn{প্রশ্ন} \bn{আছে।}
\end{quote}

\begin{quote}
\emph{ei prôsôngge āmār ekṭā prôshno āchhe.}
\end{quote}

\begin{quote}
I have a question on this topic.
\end{quote}

\begin{cousinweb}
\textbf{\bn{প্রসঙ্গ}} (\emph{prôsônggô}) is Sanskrit \emph{prasaṅga}, a connection or an occasion (Monier-Williams). In talk it is the matter in hand; \bn{এই} \bn{প্রসঙ্গে} means 'on this point'.
\end{cousinweb}

\begin{grammarlens}[title={What you've built}]
One more word for giving reasons and opinions.
\end{grammarlens}

\subsection*{Your turn}
\begin{itemize}
\raggedright
  \item \emph{Say it:} \emph{prôsônggô}
  \item \emph{Say it:} \emph{prôsônggô}, once more
  \item \emph{Say it:} \emph{ei prôsôngge āmār ekṭā prôshno āchhe.} — I have a question on this topic.
  \item \emph{Recall:} say the Bengali for an outlook, then the Bengali for an example, then say \emph{prôsônggô} again
\end{itemize}

\begin{tcolorbox}[breakable,colback=teal!4,colframe=teal!35!black,title={Before you move on}]
What does \bn{প্রসঙ্গ} mean? (\textquotedblleft{}A topic of talk; a context\textquotedblright{}.) Say it once more.
\end{tcolorbox}
